\justifying
\NSFCBodyText

以下根据申请人已确认资料列示与方法或应用背景相关的在研项目。列示用于说明研究基础和任务区别，不将既有项目成果自动计入本申请验收。

\textbf{多模态融合技术研究及其在新疆产业集群中的应用。}自治区天池英才引进计划（青年博士人才）项目，依据任务书无独立项目批准号，经费50万元，执行期为2025年2月至2028年2月；申请人主持，负责总体研究方案、关键模型设计与成果组织。该项目面向多模态技术及产业应用，本申请聚焦文本工具智能体在固定资源下的持续后训练与行动校准，拟形成独立任务协议和机制验证。

\textbf{面向政企数字化应用的智能知识库与问答系统研发。}新疆数字云创有限公司委托横向项目，编号202604140013，经费12万元，执行期为2026年6月至2026年10月；申请人主持，按合同构建知识库。合同成果和知识产权由双方共同享有。本申请以公开任务与可重置工具沙箱开展通用方法研究，不将该合同的数据、代码或知识产权未经授权移作公开资产。

\textbf{可学习乘积量化与状态空间模型协同的长上下文高效注意力机制研究。}CAAI--蚂蚁科研基金项目，编号202604140006，经费15万元，执行期为2026年2月至2027年2月；申请人参与，负责高效注意力机制研究。该项目研究模型内部注意力计算，本申请研究动态任务中的低秩适配、轨迹学习与校准，模型内部高效注意力不列为本申请新增成果。

\textbf{融合多源数据大模型驱动的光伏发电功率预测方法研究。}自治区高校基本科研业务费科研项目培育类，编号XJEDU2026P005，经费5万元，其中直接经费3.9万元、间接经费1.1万元，执行期为2026年1月至2026年12月；申请人主持，负责总体方案、模型研究与成果组织。该项目以光伏功率预测为研究终点，本申请将公开能源时序仅作为动态任务补充环境，主要终点为智能体适应、遗忘与行动风险。

本项目主要参与者及其相关在研任务待\textbf{【待补 F-PACT-02：团队名单与相关任务说明】}落实后逐项核对；同一成果的归属依据任务书、成果标注与授权材料确定。
